\chapter{نتیجه‌گیری و پیشنهاد‌ها}

\section{مقدمه}
در این فصل به جمع‌بندی پژوهش حاضر می‌پردازیم، همچنین پیشنهاداتی برای پژوهشگران آینده در این حوزه ارائه می‌دهیم.

\section{جمع‌بندی و نتیجه‌گیری}

در این پژوهش، مدل‌های یادگیری عمیق و به‌ویژه شبکه‌های عصبی گرافی برای شناسایی پیمانه‌های عملکردی در شبکه‌های تعامل پروتئین-پروتئین پیشنهاد شد. بنابر بررسی‌های انجام‌شده، تا زمان نگارش این پایان‌نامه روشی یکپارچه مبتنی بر شبکه‌های عصبی گرافی برای این مسئله ارائه نشده بود. در این تحقیق، اطلاعات زیستی پایگاه هستی‌شناسی ژن به‌صورت مؤثر و ترکیبی از دو قالب دودویی و بردارهای تعبیه به کار گرفته شد. افزون بر این موضوع، تابع هزینه پیشنهادی در این تحقیق که نسخه بهبود یافته \lr{NOCD} برای شناسایی پیمانه‌های عملکردی است، نیز تاثیر به سزایی در عملکرد این روش داشته است.

چندین مدل محاسباتی شامل شبکه عصبی گرافی پیچشی، شبکه عصبی گرافی توجه‌محور، دو نسخه بهبودیافته‌ی مدل توجه‌محور یک مدل مجهز به سازوکار پرش دانش طراحی و بررسی شدند که در میان آن‌ها، مدل توجه‌محور بهترین عملکرد را نشان داد.

برای ارزیابی، نتایج مدل‌ پیشنهادی بر روی مجموعه‌داده‌های معیار با چندین روش پیشین و نیز بهترین روش موجود در این حوزه مقایسه شد و برتری روش ارائه‌شده در تمامی معیارهای اصلی به اثبات رسید.

\section{پیشنهاد‌ها}
همان‌طور که اشاره شد در پژوهش حاضر به شناسایی پیمانه‌های عملکردی به کمک شبکه‌های عصبی گرافی و استفاده از ویژگی‌های زیستی موجود پرداختیم و روشی یکپارچه ارائه دادیم. در ادامه چندین موضوع برای ادامه پژوهش حاضر ارائه می‌شود.

در این پژوهش، برای نمایش هر پروتئین از میانگین‌گیری تعبیه‌های عبارات  GO در هر زیر‌فضا استفاده شده است. این رویکرد اگرچه ساده به نظر می‌رسد، اما از نظر محاسباتی کارا بوده است. اما می‌تواند منجر به از دست رفتن بخشی از اطلاعات زیستی نهفته در تعبیه‌ها شود. بررسی و به‌کارگیری سایر روش‌های ادغام تعبیه‌های عبارات GO در هر زیر‌فضا، به‌عنوان ویژگی‌های بازنمایی‌کننده پروتئین‌ها، می‌تواند به‌عنوان یک مسیر پژوهشی آینده مدنظر قرار گیرد و به بهبود کیفیت نمایش پروتئین‌ها منجر شود.

همچنین تمرکز اصلی این پژوهش بر شبکه‌های تعامل پروتئین-پروتئین در گونه‌ی مخمر نان بوده است، موضوعی که در بسیاری از پژوهش‌های پیشین نیز مشاهده می‌شود. تعمیم و ارزیابی روش پیشنهادی بر روی شبکه‌های تعامل پروتئین-پروتئین مربوط به سایر گونه‌ها، می‌تواند چشم‌انداز مناسبی برای بررسی قابلیت تعمیم‌پذیری و پایداری عملکرد روش فراهم آورد.

افزون بر این، هرچند روش پیشنهادی در این پژوهش با هدف شناسایی پیمانه‌های عملکردی در شبکه‌های تعامل پروتئین-پروتئین ارائه شده است، اما انتظار می‌رود در مسائل خوشه‌بندی که مستلزم شناسایی خوشه‌های هم‌پوشان هستند نیز کارایی مناسبی داشته باشد. از این رو، آزمون روش پیشنهادی در سایر مسائل خوشه‌بندی هم‌پوشان می‌تواند به‌عنوان مسیر پژوهشی دیگری در آینده دنبال شود.